\section{CM transport and its quantitative cost}
\label{sec:transport}

This section separates the precise conclusion of the cited CM-transport theorem from the cost of a chosen cycle construction. The source theorem is used as an input. The translation estimate below is a new elementary argument, and includes the nonflat maps that occur in transport.

\subsection{C01--C02: the source theorem and the auxiliary variety}

Remark 8.5 of \cite[pp.~50--51]{OAI032} invokes \cite[Corollaries 3.4 and 6.4]{OAI001}. The latter paper was not included in the supplied archive; its official October 7, 2026 PDF was obtained and directly checked. The exact printed and PDF locators are: equations (3.15)--(3.16), p.~16; Corollary 3.4 and its proof, p.~17; and the same-dimensional caution followed by Corollary 6.4 and its complete proof, p.~33.

Let $A/\overline{\Q}$ have good reduction $A_0$ at a chosen place of residue characteristic $p$, and let $\gamma$ be a rational Hodge class of codimension $r$. Corollary 3.4 supplies a positive integer $n$, a CM abelian variety $B/\overline{\Q}$ with good reduction $B_0$, a rational Hodge class $\kappa$ on $B$, and a quasi-isogeny
\[
a:A_0^n\longrightarrow B_0,
\qquad \gamma_0=(aj)^*\kappa_0,
\]
where $j:A_0\hookrightarrow A_0^n$ is the first-factor inclusion. This equality holds in every prime-to-$p$ realization and in the stated extended crystalline realization. Put $b=aj$ and choose an integer $m>0$ such that $h=mb$ is an actual homomorphism. If $Z\in\CH^r(B)_{\Q}$ represents $\kappa$, then
\begin{equation}
z_\gamma=m^{-2r}h^*Z_0\in\CH^r(A_0)_{\Q}
\label{eq:cm-source-cycle}
\end{equation}
has all the required realizations. Here $Z_0$ denotes specialization, and $h^*$ is the Chow pullback between smooth varieties, defined by the graph and refined intersection \cite{Fulton}. It requires neither flatness nor separability. The exponent is $2r$, because multiplication by $m$ acts by $m^{2r}$ on degree-$2r$ cohomology. Rational coefficients permit $p\mid m$; rational crystalline cohomology has already inverted $p$.

The CM Hodge theorem \cite[Theorem 1.1]{OAI032} supplies the algebraic representative of $\kappa$. A representative over $\overline{\Q}$ is obtained by descending through a connected Hilbert component, and its specialization is one rational cycle in all realizations. Thus this deduction uses CM transport and the CM Hodge theorem; it does not require the separate rationality theorem of \cite{OAI001}. Surjectivity of a cycle map after tensoring with $\Q_\ell$ would not itself produce this one rational representative.

There is a material qualification to the phrase ``choose a CM point in the same family.'' In the general construction, $\dim B=n\dim A$. Proposition 3.1 of \cite{OAI001} enlarges the Hodge datum using a totally real splitting field $P$ and a representation of the form $(V\otimes P)^{\oplus s}$, with $n=s[P:\Q]$. The paragraph preceding Corollary 6.4 expressly supplies neither a same-dimensional CM lift nor a map between the original generic fibers. Accordingly, the same-family conclusion in sketch 2.6 is not supplied by this corollary. This does not refute a separate same-family lifting theorem in a suitable PEL setting: such a theorem must be applied with its own integral-model and special-point hypotheses checked.

\subsection{C03: effectivity and a degree bound for nonflat pullback}

\begin{proposition}[C03: translation pullback estimate]
\label{prop:cm-pullback}
Let $A,B$ be abelian varieties over an algebraically closed field, of dimensions $g,d$. Let $h:A\to B$ be a homomorphism, $H$ an ample integral divisor on $A$, and $L$ a symmetric ample integral divisor on $B$. Let $0\leq r\leq\min(g,d)$ and let $Z$ be an effective codimension-$r$ cycle on $B$. Set
\[
J_r=H^{g-r}(h^*L)^r,\qquad
\beta(d,r)=
\begin{cases}
\dfrac{d^d}{r^r(d-r)^{d-r}},&0<r<d,\\[4pt]
1,&r=0\text{ or }r=d.
\end{cases}
\]
A general translate of $Z$ has an effective refined pullback $Y$ of codimension $r$, possibly zero. Its degree satisfies
\begin{align}
\deg_H Y&=\deg\bigl(Z\cdot h_*H^{g-r}\bigr),
\label{eq:cm-exact-degree}\\
\deg_H Y&\leq
\beta(d,r)\binom dr\frac{J_r}{L^d}\deg_L Z.
\label{eq:cm-nef-bound}
\end{align}
The factor $\beta(d,r)$ can always be replaced by $2^d$. If $H,L$ are very ample and
\[
\Delta_h=\deg_{H\boxplus L}\Gamma_h=(H+h^*L)^g,
\]
then also
\begin{equation}
\deg_HY\leq
\beta(d,r)\frac{\binom dr}{\binom gr}
\frac{\Delta_h}{L^d}\deg_L Z.
\label{eq:cm-graph-bound}
\end{equation}
No assertion of optimality among all geometric pullback estimates is made.
\end{proposition}

\begin{proof}
First let $V$ be an integral component of $Z$ and consider the translation incidence map
\[
F:A\times V\longrightarrow B,\qquad F(x,z)=h(x)-z.
\]
If $F$ is not dominant, its general fiber is empty. If it is dominant, generic flatness gives a nonempty open subset of $B$ over which $F$ is flat. The fiber over $y$ is the scheme-theoretic inverse image $h^{-1}(V+y)$. The flatness makes intersection with the regular embedding of $y$ Tor-independent; the refined pullback is therefore represented by the effective fundamental cycle of this fiber. This is the necessary justification of effectivity for a possibly nonflat $h$. Intersect the finitely many admissible open sets for the components of $Z$. Translation preserves numerical classes, so the projection formula proves \eqref{eq:cm-exact-degree}.

Choose $e>0$ so that $eH$ is very ample, and intersect $g-r$ general members of $|eH|$ to obtain an effective complete-intersection cycle $S$ of dimension $r$ on $A$. For
\[
F_S:S\times V\longrightarrow B,\qquad (x,z)\longmapsto h(x)-z,
\]
the proper pushforward formula gives
\begin{equation}
(L^d)e^{g-r}\deg_Hh^*[V+y]
=\int_{S\times V}(F_S^*L)^d.
\label{eq:cm-degree-identity}
\end{equation}
If $F_S$ is dominant, both sides count its generic degree, including inseparable length. If it is not dominant, both sides vanish. This identity does not replace scheme-theoretic multiplicities by the number of reduced points.

On $S\times V$ put
\[
P_1=\operatorname{pr}_S^*(h^*L|_S),\qquad
P_2=\operatorname{pr}_V^*(L|_V).
\]
Both are nef. The theorem of the cube for the symmetric line bundle $L$ gives
\[
F_S^*L+(h(x)+z)^*L=2P_1+2P_2;
\]
see \cite[Corollaries 6.4--6.6]{MilneAV}. Thus $2(P_1+P_2)-F_S^*L$ is nef. Differences of powers are sums of intersections of nef divisors, so
\begin{align*}
\int_{S\times V}(F_S^*L)^d
&\leq 2^d\int_{S\times V}(P_1+P_2)^d\\
&=2^d\binom dr e^{g-r}J_r\deg_LV.
\end{align*}
Only the term with exponent $r$ on $P_1$ contributes, for dimensional reasons. Combining with \eqref{eq:cm-degree-identity} and summing over components proves the bound with $2^d$.

For positive integers $u,v$, the same quadratic identity gives the sharper nef difference
\[
\left(1+\frac uv\right)P_1+
\left(1+\frac vu\right)P_2-F_S^*L
=\frac1{uv}([u]h(x)+[v]z)^*L.
\]
The preceding argument therefore replaces $2^d$ by
\[
(1+s)^r(1+1/s)^{d-r},\qquad s=u/v>0.
\]
For $0<r<d$ its minimum occurs at $s=(d-r)/r$ and equals $\beta(d,r)$. For $r=0$ take $s\to\infty$, and for $r=d$ take $s\to0$. The limiting inequalities hold because the quantity being bounded is fixed. This proves \eqref{eq:cm-nef-bound}.

Finally, all mixed intersections in
\[
\Delta_h=\sum_{i=0}^g\binom giH^{g-i}(h^*L)^i
\]
are nonnegative. Hence $\Delta_h\geq\binom grJ_r$, proving \eqref{eq:cm-graph-bound}.
\end{proof}

The graph estimate has thus been derived from a group-law degree identity. An ordinary reduced-component B\'ezout estimate alone would not justify refined multiplicities or preserve effectivity under the original nonflat pullback.

\subsection{C04--C05: a quantitative transport certificate}

\begin{corollary}[C04]
\label{cor:cm-certificate}
Use the transport data above, and fix $H$ on $A_0$ and symmetric ample $L$ on $B_0$. Suppose
\[
q\kappa=[Z^+]-[Z^-],\qquad q\geq1,
\]
where $Z^\pm$ are effective and their specialized $L$-degrees are at most $D$. With $h=mb$ and
\[
K_h=\beta(d,r)\binom dr\frac{H^{g-r}(h^*L)^r}{L^d},
\]
one has an effective-difference representative of $N\gamma_0$ with
\[
N=qm^{2r},\qquad \deg_HY^+,\deg_HY^-\leq K_hD.
\]
Consequently
\begin{equation}
c_H(\gamma_0)\leq\max\{qm^{2r},\lceil K_hD\rceil\}.
\label{eq:cm-complexity}
\end{equation}
The exact numerator degrees are
\[
D_h^\pm=\deg\bigl(Z_0^\pm\cdot h_*H^{g-r}\bigr).
\]
\end{corollary}

\begin{proof}
Descend the finitely many cycles to a number field and extend the selected place so that the abelian variety and polarization have good models. Cycle specialization preserves effectivity and the intersection degree. Apply Proposition~\ref{prop:cm-pullback} to the specialized cycles, using translations good for every component. Translation preserves their classes in all the realizations under consideration. Thus
\[
[Y^+]-[Y^-]
=q h^*\kappa_0
=qm^{2r}b^*\kappa_0
=qm^{2r}\gamma_0.
\]
The asserted degrees follow from the proposition. The translated cycles need not be rationally equivalent to the untranslated refined pullbacks; their homological classes are exactly what the complexity definition requires.
\end{proof}

The source normalization contributes another visible denominator. In \cite[equations (3.15)--(3.16), p.~16]{OAI001}, writing $d=\dim B$, the common rational similitude and the transported class are
\[
\lambda=\frac{a^*\theta_B\cdot\theta_*^{d-1}}{\theta_*^d},
\qquad \kappa=\frac{\lambda^{-r}}{[P:\Q]}\eta_B.
\]
The rational divisor classes in the ratio are ample, so $\lambda>0$. Write $\lambda=u/v$ in lowest terms and $e=[P:\Q]$. If $q\eta_B=[Z^+]-[Z^-]$ with numerator degrees at most $D$, then
\[
qeu^r\kappa=v^r([Z^+]-[Z^-]).
\]
The resulting explicit sufficient certificate is
\begin{equation}
N=qeu^rm^{2r},\qquad
\deg_HY^\pm\leq v^rK_hD.
\label{eq:cm-normalization}
\end{equation}
These are bounds before cancellation, not minimality statements. They display the quantities that must actually be controlled: a CM cycle representative, a homomorphism denominator, a polarization normalization, and the geometric size of the integral homomorphism.

\subsection{C06: noncanonicity and meaningful minimization}

\begin{proposition}[C06]
\label{prop:cm-noncanonical}
Assume $r>0$ and $\gamma_0\ne0$. A valid transport pair $(b,\kappa)$ does not determine an intrinsic ``denominator at $p$.'' Even the least integral denominator of the selected homomorphism can be made arbitrarily large at one fixed reduction place, without changing $\gamma_0$.
\end{proposition}

\begin{proof}
For every integer $M>0$, set
\[
b_M=b/M,\qquad \kappa_M=M^{2r}\kappa.
\]
Then $b_M^*\kappa_M=\gamma_0$. Since $r>0$ and $\gamma_0\ne0$, the homomorphism $b$ is nonzero. View $\Hom(A_0,B_0)$ as a lattice in its rationalization and choose a nonzero rational coordinate of $b$. Along prime integers $M$ avoiding its numerator and denominator, the least integral denominator of $b/M$ is divisible by $M$. Choosing $M^{2r}Z$ as the CM representative also makes its displayed degree cost unbounded.

There is a second ambiguity even with $(b,\kappa)$ fixed: replace a denominator-clearing integer $m$ by $Mm$. Then the integral homomorphism becomes $[M]h$. The multiplier $qm^{2r}$ and each exact generic translated-pullback degree are multiplied by $M^{2r}$. This follows from the cohomological multiplication formula and \eqref{eq:cm-exact-degree}, including when $p\mid M$.
\end{proof}

Thus one must fix a construction and its normalizations, or specify a constrained infimum over admissible data, before discussing a numerical transport cost. For example, one can minimize the right side of \eqref{eq:cm-complexity} over a prescribed collection of CM auxiliaries, polarizations, source cycle representatives and least denominator-clearing integers. Such a minimum is an upper bound for intrinsic complexity; no assertion identifies it with the intrinsic minimum over all local cycles.

The normalized rule $m^{-2r}(mb)^*$ is independent of $m$ on degree-$2r$ cohomology and on cycles modulo homological equivalence. It need not be independent of $m$ on the full Chow group, because multiplication has distinct eigenvalues on nonzero Beauville gradings. Only the required homological representative is asserted here. For a fixed arithmetic class and polarization, the eventual constancy proved in B03 concerns the intrinsic minimum, regardless of the arbitrary growth of selected transport certificates.

\subsection{C07: ordinary canonical lifts remove the marking denominator}

\begin{proposition}[C07]
\label{prop:ordinary-transport}
Let $A_0/\overline{\F}_p$ be ordinary. Assume the rational Hodge conjecture for CM abelian varieties. Every rational algebraic cycle class on $A_0$ is the specialization of a rational Hodge class, and of a rational algebraic cycle, on its canonical Serre--Tate lift $B$. The same representative works in all prime-to-$p$ and crystalline realizations. In particular, after C01 has supplied algebraicity of $\gamma_0$, ordinary transport can be chosen with $b=\id$ and $m=1$.
\end{proposition}

\begin{proof}
Descend $A_0$ to a finite field $\F_q$, enlarging it when needed. The canonical lift is functorial and lifts every endomorphism \cite[\S3]{DeligneOrdinary}. A maximal commutative semisimple subalgebra of dimension $2\dim A_0$ therefore lifts, making $B$ a CM abelian variety; it can be viewed over $\overline{\Q}$ at the selected place. Choose the corresponding complex realization. Let $g=\dim A_0$ and let $\pi$ denote Frobenius and its lifted endomorphism.

Choose $M$ divisible by the orders of all roots of unity appearing among the eigenvalues of $\pi^*/q^r$ in degree $2r$ and $\pi^*/q^{g-r}$ in degree $2g-2r$. Put
\[
V_r=\ker\bigl((\pi^*)^M-q^{Mr}\bigr)\subset H^{2r}(B(\C),\Q),
\quad
V_{g-r}=\ker\bigl((\pi^*)^M-q^{M(g-r)}\bigr).
\]
After scalar extension these are the geometric Tate subspaces on the special fiber. Frobenius is semisimple on $H^1$ and its exterior powers. Poincar\'e duality pairs an eigenvalue $\alpha$ in degree $2r$ with $q^g/\alpha$ in complementary degree, so the pairing between these two spaces is perfect.

Every element of $V_r$ is of Hodge type $(r,r)$. Indeed, the canonical-lift description \cite[\S8]{DeligneOrdinary} identifies the $H^{1,0}$ eigenlines with positive Frobenius valuation and the $H^{0,1}$ eigenlines with valuation zero. For an ordinary variety these valuations are respectively $v(q)$ and zero. An eigenmonomial of degree $2r$ whose eigenvalue is $q^r$ times a root of unity must therefore contain exactly $r$ factors of each type. Repeated eigenvalues cause no problem, since their ordinary slope determines the same Hodge type. The complementary space is treated identically.

Choose rational bases $\alpha_i$ of $V_r$ and $\beta_j$ of $V_{g-r}$. Their pairing matrix $G_{ij}$ is rational and invertible. By CM Hodge, both bases have rational algebraic representatives on $B$. Descend them to $\overline{\Q}$ through connected Hilbert parameter spaces and specialize. Their pairings remain $G_{ij}$.

Now let $z$ be a rational algebraic cycle on $A_0$ of codimension $r$. Its class is geometric Tate, hence in every realization is a linear combination of the specialized $\alpha_i$. Pair it with the specialized $\beta_j$. The answers are the same rational intersection numbers
\[
\deg\bigl(z\cdot\operatorname{sp}(\beta_j)\bigr)
\]
in every realization. Inverting the one rational matrix $G$ gives a single rational coefficient vector $(c_i)$ in every realization. Therefore $\kappa=\sum_i c_i\alpha_i$ has a rational algebraic representative whose specialization is the class of $z$ everywhere.

For $\gamma_0$, use the local algebraic cycle supplied by C01 and take the canonical lift with its identity marking. This sets $b=\id$ and $m=1$.
\end{proof}

This argument does not assume that an arbitrary $\Q_\ell$-linear Tate class is rational in canonical-lift Betti cohomology. Rationality comes from intersection with an algebraic dual basis. It also does not lift the effectiveness of a given local cycle. The degree and denominator of the characteristic-zero cycle remain uncontrolled, and the canonical CM variety and its representative can vary with the prime. Thus C07 improves the transport marking without supplying a uniform complexity bound for an arbitrary family or an unspecified class.

\subsection{C18: a normalized constant transport example for T9}

\begin{proposition}[C18]
\label{prop:T9-constant}
For the product $A=E^4$ of Section~\ref{sec:local-weil}, the same fixed CM point and identity marking give transport certificates
\[
(N,D^+,D^-)=(1,16,16)\quad\hbox{for }t_*,
\qquad (N,D^+,D^-)=(2,16,16)\quad\hbox{for }t_0=t_*/2
\]
at every good geometric reduction. Their maximum cost is the intrinsic minimum sixteen in both cases. Thus these are exact $O(1)$ transport costs as the residue prime varies.
\end{proposition}

\begin{proof}
Take the fixed CM point $x_0$ defined by this product and its product principal polarization. After the finite extension defining $\omega$ and the graph cycles, use $x_p=x_0$ and the identity marking at every good geometric reduction. The two fixed effective cycles
\[
Z^+=\Gamma_1+\Gamma_{-1},\qquad
Z^-=\Gamma_\omega+\Gamma_{-\omega}
\]
have degree sixteen each and represent $t_*$ with multiplier one. Hence the actual transport data may be chosen as $B=A$, $b=h=\id$, $m=1$, and $\kappa=t_*$. The same two cycles represent $2t_0$, yielding the second certificate. Theorem~\ref{thm:P13} gives the matching lower bound for $t_*$, and V06 gives it for $t_0$ at every good place, including characteristics two and three.
\end{proof}

No varying canonical lift is needed in this example. An arbitrary chosen arithmetic abelian variety and class require their own normalized transport data; the freedom exhibited in C06 remains present there.
